\subsection{Generators of the original integral cusp lattice}
\label{shared-actual-integral-cusp-frame}

The short theta sections span over the coefficient field.  The
arithmetic section argument also uses a finite family generating the
whole integral lattice~\cite{bergman-integral-frame-source}.
Its index in the lattice spanned by the
known sections can be arbitrarily large.  The following construction
retains that index without putting it into the norm bound.

\begin{lemma}[Integral generators from a short spanning family]
\label{lem:actual-integral-generator-conversion}
Let $K$ be a number field of degree $a$, let $V$ have dimension
$q\ge1$ over $K$, and let $\mathcal M\subset V$ be a full,
finitely generated torsion-free $\mathcal O_K$-module.  Give every
$V\otimes_{K,\sigma}\mathbb C$ a Hermitian norm.  Suppose
$\mathcal M$ contains a $K$-spanning family with norm at most
$R\ge1$ at every embedding.  Then $\mathcal M$ has at most $2aq$
generators over $\mathbb Z$, each with every embedding norm at most
\begin{equation}\label{eq:actual-integral-generator-bound}
                      a^3q\operatorname{rd}(K)R.
\end{equation}
No freeness over $\mathcal O_K$ or bound for a lattice index is
assumed.
\end{lemma}

\begin{proof}
Use the canonical Euclidean norm
$\|x\|^2=\sum_{\sigma:K\hookrightarrow\mathbb C}|\sigma x|^2$,
where both members of a conjugate pair are included.  The Euclidean
dual of $\mathcal O_K$ is the conjugate of the inverse different.
For $0\ne b\in\mathfrak D_K^{-1}$, its ideal norm gives
$|N_{K/\mathbb Q}b|\ge|\operatorname{disc}K|^{-1}$.
Arithmetic--geometric mean therefore gives dual minimum at least
$\sqrt a/\operatorname{rd}(K)$.
The elementary transference bound
$\lambda_a(L)\lambda_1(L^*)\le a^2$ gives integral elements
$b_1,\ldots,b_a$, independent over $\mathbb Q$, such that
\[
           |\sigma b_j|\le a^{3/2}\operatorname{rd}(K)
                         \le a^2\operatorname{rd}(K)
                      \quad\text{for every }\sigma,j.
\]
Select a $K$-basis $s_1,\ldots,s_q$ from the given short family.
The $aq$ elements $b_js_i$ form a $\mathbb Q$-basis and belong to
$\mathcal M$.  Their integral span $\mathcal N$ has finite index
in $\mathcal M$.

The finite abelian group $\mathcal M/\mathcal N$ has at most
$aq$ generators: the images of any integral basis of the rank-$aq$
abelian group $\mathcal M$ suffice.  Reduce a representative of
each chosen generator in the parallelepiped of the basis $b_js_i$,
so that all its real coordinates lie in $[-1/2,1/2]$.
The reduced representative remains in $\mathcal M$ and has norm
at most
\[
          \frac{aq}{2}\,a^2\operatorname{rd}(K)R
                          \le a^3q\operatorname{rd}(K)R
\]
at every embedding.  The basis vectors of $\mathcal N$ satisfy
the same bound.  These two finite families generate
$\mathcal M$ over $\mathbb Z$ and have at most $2aq$ members.
The reduction is purely in the underlying real vector space;
it subtracts integral multiples of elements of $\mathcal N$
and so introduces no denominators into $\mathcal M$.
\end{proof}

\begin{theorem}[The prescribed integral cusp frame]
\label{thm:actual-integral-cusp-frame}
Use the original good cover and coefficient field $K$, with all
the hypotheses of Theorem~\ref{thm:actual-upper-bergman} for the
first divided weight $u_1$.  Put
\[
 X=frL_b,\qquad Y=\kappa/T+1,
 \qquad T\ge1,\quad\kappa\ge0.
\]
Suppose the original numerical conditions give
\begin{gather}
 X\ge8,\qquad P\ge X^{12},\qquad
 \log[K:\mathbb Q]\le391X^2,\qquad
 \log\operatorname{rd}(K)\le391X^2,\label{eq:actual-cusp-frame-size}\\
 \sum_{\tau,i}|u_{1,\tau,i}|\le X^2Y,
 \label{eq:actual-cusp-frame-weight}
\end{gather}
and give the antitonicity and central gap $8fr+10$ required by
Theorem~\ref{thm:actual-shifted-short-span}.
Then the original integral cusp lattice on each component admits
a finite family of integral generators $w_1,\ldots,w_n$ and
numbers $\Lambda,\beta\ge1$ with
\begin{gather}
 \log n\le P,\qquad \log\beta\le P,\qquad
 \log\Lambda\le P(\kappa/T+1),
 \label{eq:actual-cusp-frame-budget}\\
 \|w_i\|_{2,Z,\sigma}\le\Lambda\quad\text{for all }i,\sigma,
 \qquad B_{Z,\sigma}(x)\le\beta.
\end{gather}
Here the norms are the original integral cusp-lattice Hermitian
norms, as identified with the component Petersson norms by the
constructed dictionary.  The coordinate matrix of these generators,
together with that Hermitian Gram matrix, is a frame for the
original lattice.  Its integral span, even over $\mathbb Z$, is
the whole lattice.

If a component's section space is zero, its integral lattice is
zero and the assertion is interpreted using one zero generator;
then $n=\Lambda=1$ is allowed.  The positive-dimensional spaces
in the original section comparison are covered without this
convention.
\end{theorem}

\begin{proof}
For a nonzero section space write $q$ for its dimension over $K$
and $a=[K:\mathbb Q]$.  The integral lattice is full and finite:
the integral model is projective, the cusp condition is the
prescribed linear restriction, and its ambient section module is
torsion-free.  These are the same finite modules used in the
short-section construction; no saturation at a finite prime is
changed here.

Theorem~\ref{thm:actual-shifted-short-span} supplies a finite
integral $K$-spanning family with all its original Hermitian norms
at most
\[
 R=\exp\left(400000\left[X^3
                   \sum_{\tau,i}|u_{1,\tau,i}|+X^5\right]\right).
\]
Apply Lemma~\ref{lem:actual-integral-generator-conversion} to
this very lattice and these norms.  Theorem
\ref{thm:actual-upper-bergman} gives
\[
           q\le\beta,\qquad \beta=\exp(4000X^3).
\]
We may therefore choose $n\le2aq$ and
$\Lambda=a^3q\operatorname{rd}(K)R$.
Writing the generators in any $K$-basis gives their required
coordinate matrix; the Gram matrix is that of the original
Hermitian norm in this basis.  The isometry and generation
properties follow from these definitions.  This construction
does not require a free integral basis.

Since $Y\ge1$,
\begin{align*}
 \log n&\le\log2+391X^2+4000X^3\le5000X^3,\\
 \log\Lambda
 &\le1564X^2+4000X^3
                      +400000X^5(Y+1)\\
 &\le806000X^5Y.
\end{align*}
For $X\ge8$ one has $5000X^3\le X^{12}$,
$4000X^3\le X^{12}$ and $806000X^5\le X^{12}$.
Together with $P\ge X^{12}$, these are precisely
\eqref{eq:actual-cusp-frame-budget}.
\end{proof}

In the original section construction, $P$ is its prescribed size
parameter and $T$ its prescribed integer length.  The existing
weight and size calculations give
\eqref{eq:actual-cusp-frame-size}--\eqref{eq:actual-cusp-frame-weight}
for the actual anchored division of the weight.  Thus the theorem
supplies the upper cusp frame required there, without a new
coefficient field, a free-lattice hypothesis or a hidden index
factor.  The second divided weight already has its finite
Hermitian lattice by the constructed positivity and section
theorems, and requires no additional quantitative frame bound.
The thin-region lower norm, point heights and final section
remainder remain open.
